\documentclass[aspectratio=169,11pt]{beamer}

\usepackage[utf8]{inputenc}
\usepackage[T1]{fontenc}
\usepackage[spanish]{babel}
\usepackage{lmodern}
\usepackage{amsmath, amssymb}
\usepackage{booktabs}
\usepackage{array}
\usepackage{tikz}
\usepackage{xcolor}
\usepackage{listings}
\usepackage{hyperref}

\usetheme{Madrid}
\usecolortheme{seahorse}
\setbeamertemplate{navigation symbols}{}
\setbeamertemplate{footline}[frame number]

\definecolor{macropolblue}{RGB}{24,68,120}
\definecolor{macropolgreen}{RGB}{56,128,98}
\definecolor{codegray}{RGB}{245,245,245}
\definecolor{darkgray}{RGB}{70,70,70}

\setbeamercolor{title}{fg=macropolblue}
\setbeamercolor{frametitle}{fg=macropolblue}
\setbeamercolor{structure}{fg=macropolblue}
\setbeamercolor{block title}{bg=macropolblue!85,fg=white}
\setbeamercolor{block body}{bg=macropolblue!5,fg=black}

\lstdefinestyle{Rstyle}{
  language=R,
  backgroundcolor=\color{codegray},
  basicstyle=\ttfamily\scriptsize,
  keywordstyle=\color{macropolblue}\bfseries,
  commentstyle=\color{darkgray}\itshape,
  stringstyle=\color{macropolgreen},
  showstringspaces=false,
  breaklines=true,
  frame=single,
  rulecolor=\color{gray!30}
}
\lstset{style=Rstyle}

\title[Sesión 2: R, datos y visualizacion]{Sesión 2: R, estructura del proyecto, limpieza de datos y visualizacion}
\subtitle{Curso: Modelos CGE y Microsimulación aplicados al Sistema de Pensiones de Capitalización Individual}
\author{MACROPOL: Economía Aplicada y Políticas Públicas}
\date{ }

\begin{document}

% 1
\begin{frame}
  \titlepage
\end{frame}

% 2
\begin{frame}{Objetivo de la sesión}
\begin{block}{Meta práctica}
Al finalizar la sesión, cada participante debe tener un proyecto en R organizado, reproducible y capaz de cargar, limpiar, transformar y visualizar una base preliminar para análisis previsional.
\end{block}

\vspace{0.3cm}
La sesión conecta tres piezas:
\begin{enumerate}
  \item \textbf{Estructura del proyecto}: cómo organizar archivos, datos y scripts.
  \item \textbf{Limpieza de datos}: cómo pasar de datos crudos a datos analíticos.
  \item \textbf{Visualizacion}: cómo producir evidencia clara para diagnóstico y política pública.
\end{enumerate}
\end{frame}

% 3
\begin{frame}{Dónde estamos en el roadmap del curso}
\centering
\begin{tikzpicture}[node distance=0.9cm, every node/.style={font=\small}]
\node[draw, rounded corners, fill=macropolblue!10, minimum width=5.0cm, minimum height=0.6cm] (A) {A. Motivación y evidencia internacional};
\node[draw, rounded corners, fill=macropolgreen!25, below of=A, minimum width=5.0cm, minimum height=0.6cm] (B) {B. Herramientas cuantitativas en R};
\node[draw, rounded corners, fill=gray!10, below of=B, minimum width=5.0cm, minimum height=0.6cm] (C) {C. Microeconomía para CGE};
\node[draw, rounded corners, fill=gray!10, below of=C, minimum width=5.0cm, minimum height=0.6cm] (D) {D. SAM y modelo CGE};
\node[draw, rounded corners, fill=gray!10, below of=D, minimum width=5.0cm, minimum height=0.6cm] (E) {E. Microsimulación previsional};
\node[draw, rounded corners, fill=gray!10, below of=E, minimum width=5.0cm, minimum height=0.6cm] (F) {F. Integración y proyecto final};
\draw[->, thick] (A) -- (B);
\draw[->, thick] (B) -- (C);
\draw[->, thick] (C) -- (D);
\draw[->, thick] (D) -- (E);
\draw[->, thick] (E) -- (F);
\end{tikzpicture}
\end{frame}

% 4
\begin{frame}{Por qué iniciar con herramientas}
\begin{block}{Idea central}
Un modelo CGE o una microsimulación no comienzan con ecuaciones. Comienzan con un flujo de trabajo ordenado.
\end{block}

\vspace{0.3cm}
Errores comunes en proyectos cuantitativos:
\begin{itemize}
  \item archivos duplicados o sin versión;
  \item bases modificadas manualmente;
  \item scripts que dependen de rutas locales;
  \item resultados que no pueden reproducirse;
  \item gráficos sin narrativa de política pública.
\end{itemize}
\end{frame}

% 5
\begin{frame}{Principio de trabajo del curso}
\centering
\Large
\textbf{Toda simulación debe ser reproducible.}

\vspace{0.6cm}
\normalsize
\begin{tikzpicture}[node distance=1.6cm, every node/.style={font=\small}]
\node[draw, rounded corners, fill=macropolblue!10, minimum width=2.1cm, minimum height=0.8cm] (data) {Datos};
\node[draw, rounded corners, fill=macropolblue!10, right of=data, xshift=1.0cm, minimum width=2.1cm, minimum height=0.8cm] (code) {Código};
\node[draw, rounded corners, fill=macropolblue!10, right of=code, xshift=1.0cm, minimum width=2.1cm, minimum height=0.8cm] (out) {Resultados};
\node[draw, rounded corners, fill=macropolblue!10, right of=out, xshift=1.0cm, minimum width=2.1cm, minimum height=0.8cm] (brief) {Decisión};
\draw[->, thick] (data) -- (code);
\draw[->, thick] (code) -- (out);
\draw[->, thick] (out) -- (brief);
\end{tikzpicture}

\vspace{0.5cm}
\normalsize
Si una persona no puede reconstruir el resultado desde los datos y el código, el resultado no está listo para política pública.
\end{frame}

% 6
\begin{frame}[fragile]{Estructura del proyecto}
\begin{lstlisting}[language=bash]
Curso_CGE_Pensiones/
|-- data/
|   |-- raw/          # datos originales, nunca se editan
|   |-- intermediate/ # datos parcialmente procesados
|   |-- clean/        # base analitica final
|
|-- R/
|   |-- 00_setup.R
|   |-- 01_importar_datos.R
|   |-- 02_limpieza_datos.R
|   |-- 03_visualizacion.R
|   |-- 04_indicadores_previsionales.R
|
|-- outputs/
|   |-- figures/
|   |-- tables/
|
|-- docs/
|-- README.md
|-- Curso_CGE_Pensiones.Rproj
\end{lstlisting}
\end{frame}

% 7
\begin{frame}{Regla de oro para los datos}
\begin{block}{Los datos originales no se modifican}
La carpeta \texttt{data/raw/} debe contener la base tal como fue recibida. Todas las transformaciones se hacen mediante código.
\end{block}

\vspace{0.3cm}
Esto permite:
\begin{itemize}
  \item auditar cambios;
  \item corregir errores sin reconstruir manualmente;
  \item actualizar bases cuando llegan nuevas rondas de información;
  \item documentar claramente el puente entre microdatos y resultados.
\end{itemize}
\end{frame}

% 8
\begin{frame}[fragile]{Script 00: configuración}
\begin{lstlisting}
# 00_setup.R
rm(list = ls())

paquetes <- c(
  "tidyverse", "data.table", "janitor", "here",
  "readxl", "scales", "skimr"
)

instalar <- paquetes[!paquetes %in% installed.packages()[, "Package"]]
if (length(instalar) > 0) install.packages(instalar)

invisible(lapply(paquetes, library, character.only = TRUE))

# Rutas principales
path_raw   <- here::here("data", "raw")
path_clean <- here::here("data", "clean")
path_fig   <- here::here("outputs", "figures")
path_tab   <- here::here("outputs", "tables")
\end{lstlisting}
\end{frame}

% 9
\begin{frame}{Datos para la sesión}
Usaremos una base sintética con información individual inspirada en microdatos laborales y previsionales.

\vspace{0.3cm}
Variables principales:
\begin{itemize}
  \item edad;
  \item sexo;
  \item nivel educativo;
  \item ingreso laboral;
  \item condición de formalidad;
  \item afiliación previsional;
  \item densidad de cotización;
  \item saldo acumulado;
  \item factor de expansión.
\end{itemize}

\vspace{0.2cm}
El objetivo no es estimar todavía el modelo final, sino aprender el flujo de trabajo.
\end{frame}

% 10
\begin{frame}[fragile]{Script 01: importar datos}
\begin{lstlisting}
# 01_importar_datos.R
source(here::here("R", "00_setup.R"))

personas_raw <- readr::read_csv(
  here::here("data", "raw", "personas_pensiones_sintetica.csv"),
  show_col_types = FALSE
)

# Estandarizar nombres de variables
personas <- personas_raw |>
  janitor::clean_names()

glimpse(personas)
skimr::skim(personas)
\end{lstlisting}
\end{frame}

% 11
\begin{frame}{Diagnóstico inicial de calidad}
Antes de limpiar, observamos.

\vspace{0.3cm}
Preguntas básicas:
\begin{itemize}
  \item ¿Cuántas observaciones tiene la base?
  \item ¿Qué variables son numéricas, categóricas o texto?
  \item ¿Hay valores faltantes?
  \item ¿Existen edades, ingresos o saldos imposibles?
  \item ¿La variable de expansión está disponible?
\end{itemize}

\vspace{0.2cm}
En pensiones, los errores de datos pueden traducirse en diagnósticos equivocados de suficiencia y cobertura.
\end{frame}

% 12
\begin{frame}[fragile]{Script 02: reglas de limpieza}
\begin{lstlisting}
# 02_limpieza_datos.R
source(here::here("R", "00_setup.R"))
source(here::here("R", "01_importar_datos.R"))

personas_clean <- personas |>
  filter(edad >= 15, edad <= 100) |>
  mutate(
    ingreso_laboral = if_else(ingreso_laboral < 0, NA_real_, ingreso_laboral),
    saldo_afp       = if_else(saldo_afp < 0, NA_real_, saldo_afp),
    formal          = if_else(formal == 1, "Formal", "Informal"),
    afiliado_afp    = if_else(afiliado_afp == 1, "Afiliado", "No afiliado"),
    grupo_edad = case_when(
      edad < 30 ~ "15-29",
      edad < 45 ~ "30-44",
      edad < 60 ~ "45-59",
      TRUE      ~ "60+"
    )
  )
\end{lstlisting}
\end{frame}

% 13
\begin{frame}[fragile]{Crear variables analíticas}
\begin{lstlisting}
personas_clean <- personas_clean |>
  mutate(
    cotizante = if_else(formal == "Formal" & afiliado_afp == "Afiliado", 1, 0),
    aporte_anual_estimado = ingreso_laboral * 12 * 0.10 * cotizante,
    quintil_ingreso = ntile(ingreso_laboral, 5),
    tasa_aporte_efectiva = aporte_anual_estimado / (ingreso_laboral * 12)
  )

readr::write_csv(
  personas_clean,
  here::here("data", "clean", "personas_pensiones_clean.csv")
)
\end{lstlisting}
\end{frame}

% 14
\begin{frame}{Variables puente hacia el modelo previsional}
\begin{table}
\centering
\small
\begin{tabular}{p{3.6cm} p{4.0cm} p{4.0cm}}
\toprule
Variable & Uso en microsimulación & Uso en CGE \\
\midrule
Ingreso laboral & salario inicial por individuo & masa salarial agregada \\
Formalidad & densidad de cotización & mercado laboral formal/informal \\
Edad & cohorte y horizonte de retiro & estructura demográfica \\
Afiliación AFP & cobertura previsional & contribuciones agregadas \\
Saldo AFP & riqueza previsional inicial & ahorro financiero \\
Factor expansión & agregación poblacional & cierre macro-micro \\
\bottomrule
\end{tabular}
\end{table}
\end{frame}

% 15
\begin{frame}{Del dato individual al indicador agregado}
\begin{block}{Ejemplo}
La cobertura previsional no es solo el porcentaje de afiliados en la muestra. Debe calcularse usando factores de expansión.
\end{block}

\vspace{0.3cm}
\[
Cobertura = \frac{\sum_i w_i \cdot Afiliado_i}{\sum_i w_i}
\]

\vspace{0.2cm}
Donde \(w_i\) es el factor de expansión del individuo \(i\).
\end{frame}

% 16
\begin{frame}[fragile]{Indicadores ponderados}
\begin{lstlisting}
indicadores <- personas_clean |>
  summarise(
    poblacion = sum(factor_expansion, na.rm = TRUE),
    cobertura_afp = weighted.mean(afiliado_afp == "Afiliado",
                                  factor_expansion, na.rm = TRUE),
    formalidad = weighted.mean(formal == "Formal",
                               factor_expansion, na.rm = TRUE),
    ingreso_promedio = weighted.mean(ingreso_laboral,
                                     factor_expansion, na.rm = TRUE),
    saldo_promedio = weighted.mean(saldo_afp,
                                  factor_expansion, na.rm = TRUE)
  )

indicadores
\end{lstlisting}
\end{frame}

% 17
\begin{frame}{Visualizacion para política pública}
Un buen gráfico debe responder una pregunta concreta.

\vspace{0.3cm}
Ejemplos:
\begin{itemize}
  \item ¿La cobertura aumenta con el ingreso?
  \item ¿La formalidad varía por edad?
  \item ¿Dónde se concentran los trabajadores sin AFP?
  \item ¿Qué grupos tendrían menor tasa de reemplazo?
  \item ¿Cómo cambia el diagnóstico antes y después de una reforma?
\end{itemize}

\vspace{0.2cm}
Visualizar no es decorar. Es revelar una relación económica.
\end{frame}

% 18
\begin{frame}[fragile]{Gráfico 1: cobertura por quintil}
\begin{lstlisting}
p1 <- personas_clean |>
  group_by(quintil_ingreso) |>
  summarise(
    cobertura = weighted.mean(afiliado_afp == "Afiliado",
                              factor_expansion, na.rm = TRUE),
    .groups = "drop"
  ) |>
  ggplot(aes(x = factor(quintil_ingreso), y = cobertura)) +
  geom_col() +
  scale_y_continuous(labels = scales::percent) +
  labs(
    title = "Cobertura AFP por quintil de ingreso",
    x = "Quintil de ingreso laboral",
    y = "Cobertura"
  ) +
  theme_minimal()
\end{lstlisting}
\end{frame}

% 19
\begin{frame}[fragile]{Gráfico 2: ingreso laboral por edad y formalidad}
\begin{lstlisting}
p2 <- personas_clean |>
  filter(!is.na(ingreso_laboral)) |>
  ggplot(aes(x = edad, y = ingreso_laboral)) +
  geom_point(alpha = 0.15) +
  geom_smooth(se = FALSE) +
  facet_wrap(~ formal) +
  scale_y_continuous(labels = scales::comma) +
  labs(
    title = "Perfil edad-ingreso por formalidad",
    x = "Edad",
    y = "Ingreso laboral mensual"
  ) +
  theme_minimal()
\end{lstlisting}
\end{frame}

% 20
\begin{frame}[fragile]{Gráfico 3: saldo AFP por grupo de edad}
\begin{lstlisting}
p3 <- personas_clean |>
  filter(afiliado_afp == "Afiliado") |>
  ggplot(aes(x = grupo_edad, y = saldo_afp)) +
  geom_boxplot(outlier.alpha = 0.2) +
  scale_y_continuous(labels = scales::comma) +
  labs(
    title = "Distribucion del saldo AFP por grupo de edad",
    x = "Grupo de edad",
    y = "Saldo acumulado"
  ) +
  theme_minimal()
\end{lstlisting}
\end{frame}

% 21
\begin{frame}[fragile]{Guardar resultados}
\begin{lstlisting}
ggsave(
  filename = here::here("outputs", "figures", "cobertura_quintil.png"),
  plot = p1, width = 8, height = 5, dpi = 300
)

ggsave(
  filename = here::here("outputs", "figures", "perfil_edad_ingreso.png"),
  plot = p2, width = 8, height = 5, dpi = 300
)

ggsave(
  filename = here::here("outputs", "figures", "saldo_afp_edad.png"),
  plot = p3, width = 8, height = 5, dpi = 300
)
\end{lstlisting}
\end{frame}

% 22
\begin{frame}{Mini-laboratorio de la sesión}
\begin{enumerate}
  \item Crear la estructura del proyecto.
  \item Cargar la base sintética de personas.
  \item Limpiar nombres y valores inválidos.
  \item Crear variables previsionales básicas.
  \item Calcular cobertura, formalidad e ingreso promedio.
  \item Construir tres gráficos de diagnóstico.
  \item Guardar la base limpia y las figuras.
\end{enumerate}
\end{frame}

% 23
\begin{frame}{Conexión con las siguientes sesiones}
Lo construido hoy alimenta directamente los módulos posteriores.

\vspace{0.3cm}
\begin{table}
\centering
\small
\begin{tabular}{p{4.0cm} p{7.0cm}}
\toprule
Producto de hoy & Uso posterior \\
\midrule
Base limpia & Microsimulación de tasas de reemplazo \\
Cobertura por grupo & Diagnóstico previsional inicial \\
Masa salarial & Contribuciones agregadas al CGE \\
Formalidad & Segmentación del mercado laboral \\
Saldo AFP & Stock inicial de riqueza previsional \\
Gráficos & Insumos para policy brief \\
\bottomrule
\end{tabular}
\end{table}
\end{frame}

% 24
\begin{frame}{Buenas prácticas desde el inicio}
\begin{itemize}
  \item Nunca editar manualmente bases crudas.
  \item Usar rutas relativas con \texttt{here()}.
  \item Separar importación, limpieza, análisis y visualizacion.
  \item Guardar resultados en carpetas estandarizadas.
  \item Nombrar scripts con orden numérico.
  \item Escribir comentarios breves, pero útiles.
  \item Pensar cada gráfico como respuesta a una pregunta.
\end{itemize}
\end{frame}

% 25
\begin{frame}{Cierre}
\begin{block}{Mensaje principal}
Antes de construir un CGE o una microsimulación, necesitamos una cadena de trabajo confiable: datos ordenados, código reproducible e indicadores interpretables.
\end{block}

\vspace{0.4cm}
En la próxima sesión entraremos al repaso microeconómico que da sentido económico a las ecuaciones del modelo.
\end{frame}

\end{document}
